\section{Configuración experimental y estrategias de entrenamiento}\label{sec:configuracion_experimental}

Para asegurar la validez de los modelos predictivos y evitar el sobreajuste (\textit{overfitting}), se diseñó una estrategia de partición de datos y validación cruzada rigurosa, adaptada a la naturaleza jerárquica de los corpus de voz.

\subsection{Partición de los datos y prevención del \textit{Data Leakage}}
La unidad muestral en este estudio no es la grabación individual, sino el paciente. Dado que cada sujeto aportó múltiples muestras de audio correspondientes a distintas tareas fonatorias (vocales, frases y habla espontánea), un muestreo aleatorio simple habría provocado la contaminación del conjunto de prueba con patrones biométricos aprendidos en el conjunto de entrenamiento, un fenómeno conocido como \textit{Data Leakage}.

Para evitarlo, la partición de los datos se realizó estrictamente a nivel de paciente (agrupando por el identificador único del locutor). Se aplicó una división del tipo \textit{Hold-out}, reservando un 80\% de los sujetos para la fase de construcción del modelo (entrenamiento y validación cruzada) y el 20\% restante como conjunto de prueba externo. Adicionalmente, esta partición se ejecutó de forma estratificada utilizando una clave combinada que integraba la clase objetivo (Parkinson o Control) y el origen del corpus (NeuroVoz o PC-GITA), garantizando así el balanceo epidemiológico en ambas particiones.

\subsection{Validación Cruzada (\textit{Cross-Validation})}
Sobre el conjunto de entrenamiento (80\% de los datos), se implementó una estrategia de validación cruzada para la optimización de los hiperparámetros. Se empleó el algoritmo \textit{Stratified Group K-Fold} con $K=5$ particiones (\textit{folds}). 

Este enfoque asegura dos propiedades matemáticas fundamentales durante cada iteración:
\begin{enumerate}
    \item \textbf{Agrupación estricta:} Ninguna grabación de un mismo paciente aparece simultáneamente en el \textit{fold} de entrenamiento y en el de validación.
    \item \textbf{Estratificación:} Se mantiene la proporción original de la variable objetivo en cada partición, mitigando los sesgos por desbalanceo de clases.
\end{enumerate}

\subsection{Optimización de umbral de decisión (Estadístico de Youden)}
De forma predeterminada, los clasificadores binarios establecen el umbral de decisión en una probabilidad de 0.5. Sin embargo, en escenarios de diagnóstico clínico, este valor rara vez es el óptimo.

Para calibrar los modelos, se separó un 20\% interno del conjunto de entrenamiento exclusivo para validación. Las predicciones probabilísticas de las grabaciones se agregaron calculando la media aritmética por paciente. Sobre estas probabilidades agregadas, se aplicó la curva ROC y se calculó empíricamente el umbral óptimo maximizando el estadístico $J$ de Youden \cite{youden1950}, definido matemáticamente como:

$$J = Sensibilidad + Especificidad - 1$$

O, equivalentemente en términos de tasas:

$$J = TPR - FPR$$

Donde $TPR$ representa la tasa de verdaderos positivos (Sensibilidad) y $FPR$ la tasa de falsos positivos ($1 - Especificidad$). El valor probabilístico de corte que maximizó $J$ fue el empleado para la binarización de las predicciones en el conjunto de prueba externo.

\subsection{Enfoque Multi-Tarea y Modelado Específico por Actividad}\label{sec:multitarea}

A nivel clínico, las alteraciones neuromotoras producidas por la disartria en la Enfermedad de Parkinson no se manifiestan de manera uniforme en todas las tareas del habla. La fonación sostenida de una vocal aísla la vibración mecánica de las cuerdas vocales, mientras que la repetición de frases o el habla espontánea introducen variables complejas como la prosodia, la competencia velofaríngea y la carga cognitiva.

Por consiguiente, mezclar todas las grabaciones en un único conjunto de entrenamiento destruiría la varianza biológica propia de cada prueba. Para preservar esta integridad y evaluar el potencial diagnóstico de cada subsistema, el pipeline de entrenamiento se ejecutó bajo un enfoque multi-tarea. Se instanciaron y evaluaron de forma completamente independiente cuatro variantes para cada algoritmo de Inteligencia Artificial (KNN, XGBoost y ResNet18), segmentadas según la actividad fonatoria de entrada:

\begin{itemize}
    \item \textbf{Modelo de Vocales Sostenidas (\textit{vocal}):} Entrenamiento restringido a grabaciones de fonación pura, diseñado para identificar anomalías en la inestabilidad de la frecuencia y la amplitud (Jitter y Shimmer).
    \item \textbf{Modelo de Frases y Lectura (\textit{frase}):} Entrenamiento enfocado en tareas de lectura guiada y repetición, evaluando la agilidad articulatoria y la transición fonética.
    \item \textbf{Modelo de Habla Espontánea (\textit{espontanea}):} Entrenamiento basado en monólogos libres, capturando la degradación de los patrones discursivos naturales.
    \item \textbf{Modelo de Fusión Global (\textit{all}):} Un clasificador holístico entrenado con el conjunto total de grabaciones, con el objetivo de determinar si la combinación de todas las modalidades acústicas incrementa la robustez de la predicción frente a los modelos aislados.
\end{itemize}

Esta separación metodológica permite analizar los resultados \textbf{\autoref{ch:resultados}} no solo en términos de qué arquitectura computacional es superior, sino también de qué prueba clínica proporciona la señal acústica más discriminativa.